\documentclass[11pt,a4paper]{article}
\usepackage[T1]{fontenc}
\usepackage[utf8]{inputenc}
\usepackage{amsmath,amssymb,amsthm}
\usepackage[margin=1in]{geometry}
\usepackage[colorlinks=true,linkcolor=blue,urlcolor=blue]{hyperref}
\theoremstyle{plain}
\newtheorem{theorem}{Theorem}
\newtheorem{lemma}[theorem]{Lemma}
\newtheorem{proposition}[theorem]{Proposition}
\newtheorem{corollary}[theorem]{Corollary}
\theoremstyle{definition}
\newtheorem{remark}[theorem]{Remark}

\title{The empirical recurrence for OEIS A209118, proved from an exact model}
\author{Adrian Perez Fontelles\\ \small Independent researcher}
\date{13 September 2026}
\begin{document}
\maketitle

\begin{abstract}
OEIS A209118 carries an empirical recurrence of order $12$. It is true, and it is
decidable rather than empirical. The entry's sequence is given exactly by a Burnside sum over the necklace group, from
which $a$ provably satisfies a MONIC linear recurrence of order $S=67$, derived below and not
assumed. The residual of the conjectured recurrence satisfies the same one, so whether it
annihilates $a$ is settled by a finite exact computation. No transfer matrix is involved and
the count is not a walk.
\end{abstract}

\noindent\small 2020 Mathematics Subject Classification. 05A15, 05A19, 11P21.\normalsize

\section{The conjecture}

OEIS A209118 is ``Number of 6-bead necklaces labeled with numbers -n..n not allowing reversal, with sum zero and avoiding the patterns z z+1 z+2 and z z-1 z-2.''. It has offset $1$ and begins
\[
15, 238, 1373, 5032, 14037, 32802, 67681, 127310,\ \dots
\]
The entry states, as a conjecture contributed by its author and never marked settled:
\begin{quote}
Empirical: a(n) = 4*a(n-1) - 6*a(n-2) + 5*a(n-3) - 5*a(n-4) + 7*a(n-5) - 8*a(n-6) + 7*a(n-7) - 5*a(n-8) + 5*a(n-9) - 6*a(n-10) + 4*a(n-11) - a(n-12) for n > 13.
\end{quote}
As of the ``Last modified'' line on the live entry (Jul 08 2018 21:31:06, revision 13)
this is still recorded as empirical, and nothing on the entry records it as proved.

\section{The count is a Burnside sum}

The entry counts $6$-bead necklaces labelled from $-n..n$ with sum zero under a further
condition on consecutive beads, the beads identified under the cyclic group of rotations, written $G$. Nothing is a
walk here: the bead count is fixed and the ALPHABET grows.

By Burnside's lemma the number of orbits is $\frac{1}{|G|}\sum_{g\in G}|\mathrm{Fix}(g)|$.
A labelling fixed by $g$ is constant on the orbits of $g$, so it is a choice of one value per
orbit. For a rotation by $d$ with $c=\gcd(6,d)$ the fixed labellings are exactly the
$c$-periodic words, the whole sum is $(6/c)$ times the sum of one period, and the condition on
the $6$-periodic extension is the condition on the length-$c$ cyclic word; for a reflection
they are the palindromes, determined by $\lceil(6+1)/2\rceil$ values, and the condition is
again local in those. Every term is therefore a count of integer points of
\[
\sum_i w_i y_i = 0, \qquad -n \le y_i \le n,
\]
$w_i$ the orbit sizes, under a condition on consecutive beads. The forbidden windows the entry
names are defined by EQUALITIES between bead values, with fixed offsets, so inclusion-exclusion
over them turns each term into a count under a set of affine identifications --- one
convolution each --- and nothing is enumerated over the alphabet.

\section{The bound}

Each term is a lattice-point count over a union of relatively open rational cones in the orbit
variables, hence a quasi-polynomial in $n$; a Burnside average of quasi-polynomials is one. The
period is read off the arrangement: every ray of every cell is cut out by as many independent
hyperplanes as there are variables, so by Cramer its primitive generator has its $n$-coordinate
dividing the determinant of their $y$-parts, and a ray leaving the box or the sum-zero plane
bounds no cell and is dropped. The multiplicity of each cyclotomic factor is bounded by the
number of distinct rays whose height it divides, and by the dimension of the cone. Taking the
maximum over $g$ of the resulting multiplicities gives a monic annihilator of order
$S=67$, which includes the pre-period $n_0=18$ the conditions with an offset need. The
annihilator was then asked for terms the model had not supplied and reproduced them.

\section{The proof}

\begin{lemma}\label{lem:crit}
Let $A(z)=z^S-\sum_{j<S}\alpha_jz^{j}$ be the monic annihilator derived above and
$q(z)=z^{12}-\sum_ic_iz^{12-i}$ the characteristic polynomial of the conjectured
recurrence. The residual $r(n)=a(n)-\sum_ic_ia(n-i)$ is a linear combination of shifts of $a$,
so $A$ annihilates $r$ as well. Since $A(0)\ne0$ the recurrence it gives runs in both
directions, and $S$ consecutive zeros of $r$ force $r$ to vanish at every later index.
\end{lemma}

\begin{proof}
$A$ annihilates $a$, hence every shift of $a$, hence any linear combination of shifts; that is
$r$. A monic recurrence with nonzero constant term determines each term from the $S$ before it
and each from the $S$ after it, so a run of $S$ zeros propagates.
\end{proof}

\begin{theorem}
\[
a(n)\;=\;4\,a(n-1) - 6\,a(n-2) + 5\,a(n-3) - 5\,a(n-4) + 7\,a(n-5) - 8\,a(n-6) + 7\,a(n-7) - 5\,a(n-8) + 5\,a(n-9) - 6\,a(n-10) + 4\,a(n-11) - a(n-12)
\]
for every $n>13$.
\end{theorem}

\begin{proof}
The terms $a(0),\dots$ were computed exactly in integer arithmetic from the model, extended by
$A$ where the model itself was not evaluated, and $r(n)$ formed for each index. Those integers
vanish from $n=13+1$ onwards, and the run exceeds $S+12$, which by
Lemma~\ref{lem:crit} settles every larger index at once. The bound is exact: at $n=13$ the recurrence fails on the entry's own published terms, so no larger range holds.
\end{proof}

\section{Verification}

Three checks, each independent of the algebra above.

First, the model reproduces all $28$ terms the entry publishes, exactly, in integer
arithmetic. This is what ties the model to the entry: it is built by reading the entry's
English, and a misreading gives different counts.

Second, the conjectured recurrence was evaluated directly on those published terms, with no
model involved, and holds wherever the proved range and the data overlap.

Third, the derived annihilator $A$ was itself tested on terms it had not been given: the model
was evaluated beyond the $S$ values that determine it and $A$ reproduced them. A bound that is
too small fails this test, which is why it is run.

\begin{thebibliography}{9}
\bibitem{oeis} The OEIS Foundation, \emph{The On-Line Encyclopedia of Integer
Sequences}, \url{https://oeis.org}, 2026. Sequence A209118.
\bibitem{beckrobins} M.~Beck and S.~Robins, \emph{Computing the Continuous Discretely},
2nd ed., Springer, 2015. (Ehrhart quasi-polynomials and their periods.)
\bibitem{stanley} R.~P.~Stanley, \emph{Enumerative Combinatorics}, Volume 1, 2nd ed.,
Cambridge University Press, 2011.
\bibitem{ds} F.~Diamond and J.~Shurman, \emph{A First Course in Modular Forms},
Springer GTM 228, 2005.
\end{thebibliography}

\end{document}
